% Preamble requirements:
%   \usepackage{tikz,pifont,graphicx,textcomp}
%   \usetikzlibrary{arrows.meta,positioning,fit,calc,shapes.callouts}
% Real example, CodeGemma-7B, held-out stack C3_L1r_S1_X1, item apps_3468_0::C3_L1r_S1_X1::2
% (APPS `scramble`). Program is the entry function verbatim except: two long conditional
% expressions are line-wrapped in parentheses, the unreachable RuntimeError message is elided,
% and the unused helper defs (_rs, _ar_p, _ar_m, _ar_x) are omitted.
% The value/call in each bubble is the model's verbatim reply; the rest of the sentence is our
% explanation (the eval stores answers only, no reasoning). Sources:
%   forward : results/cells/{f2_divergence,merge_panel}/codegemma-7b/python/<arm>__C3_L1r_S1_X1
%   backward: results/cells/inverse_generic/codegemma-7b/python/<arm>__C3_L1r_S1_X1
% Middle row = cons_lam3 and mono_all (identical replies on this item). The router is NOT in the
% middle row: it answers forward correctly but backward writes f_eba3([1,2,3,3], [1,1,2,2]),
% which returns False -- a wrong call, not forward-locking (3.8% locked on this condition).
\begin{figure*}[t]
\centering
\resizebox{\textwidth}{!}{%
% ---- Colors (shared with the other obfuscation figures) ----
\definecolor{renC}{HTML}{3B6FB6}\definecolor{flowC}{HTML}{D9822B}\definecolor{exprC}{HTML}{8A5AB8}
\colorlet{renF}{renC!18}\colorlet{flowF}{flowC!22}\colorlet{exprF}{exprC!20}
\definecolor{okC}{HTML}{2E8B57}\definecolor{badC}{HTML}{C0392B}
\pgfdeclarehorizontalshading{tricolor}{100bp}{%
  color(0bp)=(renF); color(38bp)=(renF); color(50bp)=(flowF);
  color(62bp)=(exprF); color(100bp)=(exprF)}
\providecommand{\hlstrut}{\rule[-0.32em]{0pt}{1.08em}}
\providecommand{\rn}[1]{{\setlength{\fboxsep}{0.6pt}\smash{\colorbox{renF}{\hlstrut #1}}}}
\providecommand{\fl}[1]{{\setlength{\fboxsep}{0.9pt}\smash{\colorbox{flowF}{\hlstrut #1}}}}
\providecommand{\ex}[1]{{\setlength{\fboxsep}{0.6pt}\smash{\colorbox{exprF}{\hlstrut #1}}}}
\providecommand{\ind}[1]{\hspace*{#1\fontcharwd\font`\ }}
\newcommand{\ok}{{\color{okC}\ding{51}}}
\newcommand{\no}{{\color{badC}\ding{55}}}
\begin{tikzpicture}[
  font=\sffamily\small,
  title/.style={font=\sffamily\bfseries\small, anchor=base west},
  task/.style={font=\sffamily\bfseries\small, anchor=base west},
  bubble/.style={rectangle callout, rounded corners=3pt, draw=black!45, fill=white,
     align=left, font=\sffamily\footnotesize\frenchspacing, inner sep=4pt, text width=5.5cm,
     },
  bad/.style={bubble, draw=badC!70, fill=badC!4},
  good/.style={bubble, draw=okC!70, fill=okC!5},
  mname/.style={font=\sffamily\bfseries\scriptsize, align=center},
  % little model face; #1 = fill style. botgood smiles, botbad frowns
  botbase/.pic={
    \draw[black!55, thick] (0,0.42) -- (0,0.58); \fill[black!55] (0,0.62) circle (0.05);
    \draw[draw=black!60, rounded corners=3pt, #1, thick] (-0.36,-0.3) rectangle (0.36,0.42);
    \fill[white] (-0.15,0.12) circle (0.085) (0.15,0.12) circle (0.085);
    \fill[black!75] (-0.15,0.12) circle (0.04) (0.15,0.12) circle (0.04);
  },
  botgood/.pic={\pic{botbase={#1}}; \draw[black!70, thick, line cap=round] (-0.13,-0.08) to[bend right=45] (0.13,-0.08);},
  botbad/.pic={\pic{botbase={#1}}; \draw[black!70, thick, line cap=round] (-0.13,-0.16) to[bend left=45] (0.13,-0.16);},
]
% ================= Program and tasks =================
\node[draw=black!40, rounded corners=2pt, inner sep=6pt, anchor=north west] (code) at (0,0) {%
\ttfamily\scriptsize
\begin{tabular}{@{}l@{}}
def \rn{f\_eba3}(\rn{v\_9beb}, \rn{v\_9857}):\\
\ind{4}\fl{\_sn\_e705 = object()}\\
\ind{4}\fl{\_st\_bbb9 = \ex{-11891 + 11903}}\\
\ind{4}\fl{while \_st\_bbb9 != \ex{-(-20148 + 20149)}:}\\
\ind{8}\fl{if \_st\_bbb9 == \ex{44042 - (44042 - 70)}:}\\
\ind{12}return True\\
\ind{8}\fl{elif \_st\_bbb9 == \ex{23300 - (23300 - 26)}:}\\
\ind{12}\rn{v\_dbfa} = \fl{\_v\_5220}\\
\ind{12}\fl{\_st\_bbb9 = \ex{-2854 + 2917}}\\
\ind{8}\fl{elif \_st\_bbb9 == \ex{-26285 + 26348}:}\\
\ind{12}\fl{\_st\_bbb9 = (\ex{24950 - (24950 - 18)}}\\
\ind{16}\fl{if} \rn{v\_9beb}.count(\rn{v\_dbfa}) < \rn{v\_9857}.count(\rn{v\_dbfa})\\
\ind{16}\fl{else \ex{45625 - (45625 - 15)})}\\
\ind{8}\fl{elif \_st\_bbb9 == \ex{52158 - (52158 - 12)}:}\\
\ind{12}\fl{\_it\_2ee9 = iter(}set(\rn{v\_9857})\fl{)}\\
\ind{12}\fl{\_st\_bbb9 = \ex{30429 - (30429 - 15)}}\\
\ind{8}\fl{elif \_st\_bbb9 == \ex{3334 - (3334 - 18)}:}\\
\ind{12}return False\\
\ind{8}\fl{elif \_st\_bbb9 == \ex{33788 - (33788 - 15)}:}\\
\ind{12}\fl{\_v\_5220 = next(\_it\_2ee9, \_sn\_e705)}\\
\ind{12}\fl{\_st\_bbb9 = (\ex{8050 - (8050 - 26)} if \_v\_5220 is not \_sn\_e705}\\
\ind{16}\fl{else \ex{61435 - (61435 - 70)})}\\
\ind{8}\fl{else:}\\
\ind{12}\fl{raise RuntimeError(...)}
\end{tabular}};
\node[title] (ta) at ($(code.north west)+(0,0.78)$) {Stacked-obfuscated program};
\node[font=\sffamily\footnotesize, anchor=base west, inner sep=0pt] at ($(code.north west)+(0,0.3)$) {%
  {\setlength{\fboxsep}{1.5pt}\fcolorbox{renC}{renF}{\strut\texttt{L1r} renaming}}%
  \ $\rightarrow$\ %
  {\setlength{\fboxsep}{1.5pt}\fcolorbox{flowC}{flowF}{\strut\texttt{S1} flattening}}%
  \ $\rightarrow$\ %
  {\setlength{\fboxsep}{1.5pt}\fcolorbox{exprC}{exprF}{\strut\texttt{X1} arithmetic}}};
% ================= Grid of models x tasks =================
\def\xA{10.4}\def\xB{17.7}   % face x-positions for the two task columns
\def\bw{0.72}              % gap from face to bubble
\node[task] at ({\xA-0.45},0.78) {Forward reasoning};
\node[font=\sffamily\footnotesize, anchor=base west] at ({\xA-0.45},0.3) {Output prediction: \texttt{f\_eba3(s1, s2) = ?}};
\node[font=\sffamily\footnotesize, anchor=base west] at ({\xA-0.45},-0.12) {\texttt{s1=\textquotesingle cedewaraaossoqqyt\textquotesingle, s2=\textquotesingle codewars\textquotesingle}};
\node[task] at ({\xB-0.45},0.78) {Reverse reasoning};
\node[font=\sffamily\footnotesize, anchor=base west] at ({\xB-0.45},0.3) {Input prediction: find \texttt{s1, s2} such that};
\node[font=\sffamily\footnotesize, anchor=base west] at ({\xB-0.45},-0.12) {\texttt{f\_eba3(s1, s2) == True}};
\newcommand{\cell}[5]{% x, y, face style, bubble style, text
  \pic at (#1,#2) {bot#4={#3}};
  \node[#4, anchor=west, callout absolute pointer={({#1+0.42},{#2+0.05})}] at ({#1+\bw},{#2+0.05}) {#5};}
\newcommand{\rowname}[2]{\node[mname, anchor=north] at (\xA,{#1-0.38}) {#2};}

% Base model
\rowname{-1.3}{Base}
\cell{\xA}{-1.3}{fill=black!8}{bad}{Predicts \texttt{False}, but every letter of\\\texttt{codewars} occurs often enough in \texttt{s1}. \no}
\cell{\xB}{-1.3}{fill=black!8}{bad}{Calls \texttt{f\_eba3(\{1,2,3\}, \{1,2,3\})}, which\\raises: a set has no \texttt{.count}. \no}
% KL-consistency / pooled SFT
\rowname{-3.9}{KL-consistency\\Pooled SFT}
\cell{\xA}{-3.9}{fill=cyan!10}{good}{Predicts \texttt{True}: every letter of\\\texttt{codewars} occurs often enough in \texttt{s1}. \ok}
\cell{\xB}{-3.9}{fill=cyan!10}{bad}{Replies \texttt{True}, the output itself,\\instead of a call to \texttt{f\_eba3}. \no\\[1pt]{\color{badC!80!black}\itshape\scriptsize answers the trained forward task instead}}
% Merged
\rowname{-6.5}{Merged}
\cell{\xA}{-6.5}{shading=tricolor}{good}{Predicts \texttt{True}: every letter of\\\texttt{codewars} occurs often enough in \texttt{s1}. \ok}
\cell{\xB}{-6.5}{shading=tricolor}{good}{Calls \texttt{f\_eba3([1,2,3], [1,2,3])}, which\\returns \texttt{True}: every count is equal. \ok}
\end{tikzpicture}}
\caption{Qualitative example of capability preservation under model merging, taken from the
evaluation: CodeGemma-7B on a held-out program from the three-deep stack
\texttt{L1r}\,$\rightarrow$\,\texttt{S1}\,$\rightarrow$\,\texttt{X1} (APPS \texttt{scramble}:
does \texttt{s1} contain enough of each character of \texttt{s2}?). Each bubble gives the
model's verbatim answer with our explanation of it. The untuned model fails forward prediction,
while the KL-consistency and pooled-SFT arms predict the output correctly. Asked to reason in
the reverse direction, however, they reply with the output value itself rather than an input,
effectively answering the forward task they were trained on; on this stack the KL arm does so on
88\,\% of items and pooled SFT on 25\,\%. The DARE-TIES merge handles both directions.
Backward answers are graded by execution. The listing wraps two long lines and omits unused
helper definitions.}
\label{fig:case-study}
\end{figure*}
